\documentclass{category-template}

\usepackage{calrsfs}
\DeclareMathAlphabet{\pazocal}{OMS}{zplm}{m}{n}

\begin{document}

\begin{center}
    \textbf{\Huge{Sheaves and Cohomology Workbook}}

    \textit{by Dmytro Zakharov}
\end{center}

\subsubsection*{Remark}

\begin{problem}
    Denote by $\text{Mod}^f(A)$ the full subcategory of 
    $\text{Mod}(A)$ consisting of finitely generated modules.
    Does the category $\text{Mod}^f(A)$ admit kernels,
    cokernels, finite/infinite direct sums, products?
\end{problem}

\textbf{Solution.} 

\begin{itemize}
    \item \textbf{Kernels.} Kernel of the morphism $N \xrightarrow{\phi} M$
    between two finitely generated modules is not always finitely generated.
    Take $A=N=\mathbf{k}[x_1,x_2,\dots]$ and $M := \mathbf{k}[x_1,x_2,\dots]/I$
    where $I=(x_1,x_2,\dots)$. We have morphism $A \to A/I$ which has a kernel
    $I$. Surely, $I$ is not finitely generated as an $A$-module, but $A/I \cong
    \mathbf{k}$ is finitely generated.
    \item \textbf{Cokernels.} Assume $N \xrightarrow{\psi} M$ is another
    morphism of f.g. modules. Is $M/\text{im}(\psi)$ a finitely generated
    module? Yes. Note that $\text{im}(\psi)$ is finitely generated by
    $\psi(n_1),\dots,\psi(n_r)$ where $N=(n_1,\dots,n_r)$. Quotient of two
    finitely generated modules is finitely generated.
    \item \textbf{Finite direct sums.} Obviously $\bigoplus_{i \in I}N_i$ for
    finite $I$ is finitely generated by the elements
    $e_{ij}=(0,0,\dots,n_{ij},\dots,0)$ where $i \in I$ and $j$ ranges over
    generators of $N_i$.
    \item \textbf{Infinite direct sums.} For instance,
    $\mathbb{Z}^{\oplus\mathbb{N}}$ is not finitely generated.
    \item \textbf{Finite products.} Since finite products are the same as finite
    direct sums in $\text{Mod}(A)$, we have already covered this case.
    \item \textbf{Infinite products.} Take $\prod_{i \in \mathbb{N}}\mathbb{Z}$.
    It is not finitely generated.
\end{itemize}

\textbf{Conclusion:} $\text{Mod}^f(A)$ only admits cokernels, finite direct sums
and products.

\begin{problem}
    Let $X$ be a set and $\mathcal{C}$ be a category $\mathcal{P}(X)$ of all
    subsets of $X$, the morphism being the inclusion morphism. Does
    $\mathcal{C}$ admit an initial object, terminal object, products,
    coproducts?
\end{problem}

\textbf{Solution.} $P \in \mathcal{C}$ is initial if
$\text{Hom}_{\mathcal{C}}(P,X) \cong \{\mathsf{pt}\}$ for all $X \in
\mathcal{C}$. By definition recall:
\begin{equation*}
    \text{Hom}_{\mathcal{C}}(P,X) = \begin{cases}
        \{\mathsf{pt}\} & P \subseteq X, \\
        \emptyset & \text{otherwise}.
    \end{cases}
\end{equation*}

Thus $P \subseteq X$ for any $X \in \mathcal{P}(X)$. This is the case when
$P=\emptyset$.

$P$ is terminal if $\text{Hom}_{\mathcal{C}}(X,P) \cong \{\mathsf{pt}\}$ for all
$X \in \mathcal{P}(X)$. This is equivalent to $X \subseteq P$ for all $X \in
\mathcal{P}(X)$. This is the case when $P=X$.

How to define the product $X \times Y$? By definition, the following diagram
should commute:
\[
\begin{tikzcd}
& Z \arrow[dashed]{d}[description]{\exists!} \arrow[bend left]{dr}{g} \arrow[bend right]{dl}[swap]{f} & \\
X & X \times_{\mathcal{C}} Y \arrow{l}[swap]{\pi_X} \arrow{r}{\pi_Y} & Y
\end{tikzcd}
\]

So naturally $X \times_{\mathcal{C}} Y = X \cap Y$. Similarly, $X
\oplus_{\mathcal{C}} Y = X \cup Y$. Obviously, $\bigcap_{i \in I} X_i$ and
$\bigcup_{i\in I} X_i$ are always defined in $\mathcal{P}(X)$ for any
cardinality of $I$.

\begin{problem}
    Let $X$ be a topological space and let $b \in X$. One defines a presheaf $F$
    on $X$ by setting $F(U)=0$ if $b \not\in U$ and $F(U) = \mathbf{k}$ if $b
    \in U$. Prove that $F$ is a sheaf and give a short description of it by
    using the direct image functor.
\end{problem}

\textbf{Solution.} $F$ is a skyscrapper sheaf $\delta_b$ that is defined as
follows: let $i_b: \{b\} \hookrightarrow X$ be the embedding into $X$, then $F
\cong i_{b*}i_{b}^{-1}\mathbf{k}_X$. By definition, $f^{-1}F$ is a map that
associates a sheaf to a presheaf $f^{\dag}F$, thus $i_b^{-1}\mathbf{k}_X$ is a
sheaf. Since $f_*$ sends sheaf to a sheaf, the composition yields a sheaf.

\begin{problem}
    Let $X$ be an open subset of $\mathbb{C}$ and let $b \in X$. For $U$ open in
    $X$, one denotes by $\mathcal{O}_X^b(U)$ the subspace of $\mathcal{O}_X(U)$
    consisting of holomorphic functions vanishing at $b$ (Hence, if $b \not\in
    U$, $\mathcal{O}_X^b(U) = \mathcal{O}_X(U)$). Show that $U \mapsto
    \mathcal{O}_X^b(U)$ is a sheaf, that we denote by $\mathcal{O}_X^b$ and
    describe the sheaf $F^b$ given by the exact sequence of sheaves $0 \to
    \mathcal{O}_X^b \to \mathcal{O}_X \to F^b \to 0$.
\end{problem}

\textbf{Solution.} By definition $\mathcal{O}_X^b(U) = \{f \in \mathcal{O}_X(U):
f(b)=0\}$. Let us check sheaf conditions:
\begin{itemize}
    \item \textit{Locality.} Obvious.
    \item \textit{Glueing.} Suppose $\{U_i\}_{i \in U}$ is an open cover of $U$,
    and family $\{s_i \in \mathcal{O}^b_X(U_i)\}_{i \in I}$ satisfies the
    property $s_i|_{U_{ij}}=s_j|_{U_{ij}}$ for all $i,j$. By the property of
    $\mathcal{O}_X$, there exists some $s \in \mathcal{O}_X$ such that $s|_{U_i}
    = s_i$. We shall check that $s \in \mathcal{O}_X^b$, that is $s(b)=0$ (of
    course assume $b \in U$, otherwise the statement is obvious). But then there
    exists some $U_i$ such that $s|_{U_i}(b)=0$, meaning $s(b)=0$.
\end{itemize}

Now consider the exact sequence
\begin{equation*}
    0 \to \mathcal{O}_X^b \to \mathcal{O}_X \to F^b \to 0.
\end{equation*}

By the definition of the exact sequence, $F^b = \text{Coker}(\mathcal{O}_X^b \to
\mathcal{O}_X)$. Take some open $U$ and by definition of the cokernel, $F^b$ is
a sheafication of $U \mapsto \mathcal{O}_X(U) / \mathcal{O}_X^b(U)$. If $b
\not\in U$, then this corresponds to a constant sheaf $U \mapsto 0$. Otherwise,
this quotient defines the equivalence classes of functions that agree on $b$, 
thus this is isomorphic to $\mathbb{C}$. Therefore, we have the skyscrapper 
sheaf $F^b = i_{b*}i_{b}^{-1}\mathbb{C}_X$.

\begin{problem}
    Let $X=\mathbb{R}$ and consider the morphism of sheaves $P: C_X^{\infty} \to
    C_X^{\infty}$, where $P$ is the differential operator $\partial_x-1$.
    \begin{enumerate}
        \item[(i)] Calculate $\text{Ker}(P)$ and $\text{Coker}(P)$.
        \item[(ii)] Same question with $P$ replaced with $Q=x\partial_x - \frac{1}{2}$. 
    \end{enumerate}
\end{problem}

\textbf{Solution.} \textit{Part (i).} By definition, $\text{Ker}(P) = \{f \in
C_X^{\infty}: P(f) = 0\}$. Thus we obtain $f'-f=0$ or $f=\gamma e^x$ for $\gamma
\in \mathbb{R}$. Thus, $\text{Ker}(P) = \mathbb{R}e^x$. Now recall
$\text{Coker}(P) = C_X^{\infty} / \text{Im}(P)$. So what is $\text{Im}(P)$? Note
$g \in \text{Im}(P)$ iff $g = f'-f$ has a solution w.r.t $f$. But this is always
the case, so $\text{Im}(P) = C_X^{\infty}$ and so $\text{Coker}(P) = 0$.

\textit{Part (ii).} Again, $\text{Ker}(P) = \{f \in C_X^{\infty}: P(f) = 0\}$.
Thus we have $xf'(x) - \frac{1}{2}f(x) = 0$ or $f(x)=c\sqrt{|x|}$. But note that
$f(x) \not\in C_X^{\infty}$. Thus, $\text{Ker}(P) = 0$.

Now consider $\text{Im}(P)$, so we have the equation
\begin{equation*}
    x\frac{df}{dx} - \frac{1}{2}f(x) = g(x)
\end{equation*}

One can verify that:
\begin{equation*}
    f(x) = \sqrt{|x|}\left(\gamma + \int_0^x \frac{g(t)dt}{t\sqrt{|t|}}\right), \quad x \neq 0
\end{equation*}

So whenever we have the equation on some open $U \not\ni 0$, then we always have
a solution. Near $0$, however, there are some non-trivial solutions, which we
cannot describe. Thus, $\text{Coker}(P) \cong i_*\mathbb{R}_{\{0\}}$.

\begin{problem}
    Let $M,N \in \underline{\text{Mod}}(\mathbf{k})$ and consider the functor
    $\theta: \underline{\text{Mod}}(\mathbf{k}) \to
    \underline{\text{Mod}}(\mathbf{k})$ given by $L \mapsto
    \text{Hom}_{\mathbf{k}}(M \otimes_{\mathbf{k}} L, N)$. Is this functor 
    covariant/contravariant? Left/right exact?
\end{problem}

\textbf{Solution.} We shall think how to map morphisms $f \in
\text{Hom}_{\mathbf{k}}(L, R)$ to $\theta(f): \text{Hom}_{\mathbf{k}}(M
\otimes_{\mathbf{k}} R, N) \to \text{Hom}_{\mathbf{k}}(M \otimes_{\mathbf{k}} L,
N)$. Recall that 
\begin{equation*}
    \text{Hom}_{\mathbf{k}}(M \otimes_{\mathbf{k}} L, N) \cong \text{Hom}_{\mathbf{k}}(M, \text{Hom}_{\mathbf{k}}(L, N)) \cong \text{Hom}_{\mathbf{k}}(L, \text{Hom}_{\mathbf{k}}(M, N))
\end{equation*}

So for $f \in \text{Hom}_{\mathbf{k}}(L, R)$ define 
\begin{equation*}
    \theta(f): \text{Hom}_{\mathbf{k}}(R, \text{Hom}_{\mathbf{k}}(M, N)) \xrightarrow{\circ f} \text{Hom}_{\mathbf{k}}(L, \text{Hom}_{\mathbf{k}}(M, N)).
\end{equation*}

Thus $\theta$ is a contravariant functor. Now, let us assume that
\begin{equation*}
    0 \to L' \to L \to L'' \to 0
\end{equation*}
is a short exact sequence. Let us apply functor $\theta$, which is a composition of
functor $\text{Hom}_{\mathbf{k}}(\bullet, N)$ and $\text{Hom}_{\mathbf{k}}(M,
\bullet)$. Functor $\text{Hom}_{\mathbf{k}}(\bullet, N)$ is left exact 
contravariantly, thus getting the exact sequence:
\begin{equation*}
    0 \to \text{Hom}_{\mathbf{k}}(L'', N) \to \text{Hom}_{\mathbf{k}}(L, N) \to \text{Hom}_{\mathbf{k}}(L', N)
\end{equation*}

The functor $\text{Hom}_{\mathbf{k}}(M, \bullet)$ is right adjoint to the 
functor $(\bullet) \otimes_{\mathbf{k}} M$, thus is left exact, producing the short exact sequence:
\begin{equation*}
    0 \to \text{Hom}_{\mathbf{k}}(M, \text{Hom}_{\mathbf{k}}(L'', N)) \to \text{Hom}_{\mathbf{k}}(M, \text{Hom}_{\mathbf{k}}(L, N)) \to \text{Hom}_{\mathbf{k}}(M, \text{Hom}_{\mathbf{k}}(L', N)),
\end{equation*}
which is isomorphic to
\begin{equation*}
    0 \to \text{Hom}_{\mathbf{k}}(M \otimes L'', N) \to \text{Hom}_{\mathbf{k}}(M \otimes L, N) \to \text{Hom}_{\mathbf{k}}(M \otimes L', N).
\end{equation*}
Thus functor $\theta$ is left exact contravariantly.

\begin{problem}
Give an example of an additive category which admits neither kernels nor
cokernels
\end{problem}

\textbf{Solution.} Define $\mathcal{C}$ such that $\text{Ob}(\mathcal{C}) =
\mathbb{Z}$ and $\text{Mor}(\mathcal{C})$ are defined as:
\begin{equation*}
    \text{Hom}_{\mathcal{C}}(n, m) = \begin{cases}
        \mathbb{Z}, & m=n, \\
        \{0\}, & m \neq n.
    \end{cases}
\end{equation*}

Note that $\text{Mor}(\mathcal{C}) \in \underline{\text{Mod}}(\mathbb{Z})$,
composition is bilinear, and finite biproduct exist. Thus, $\mathcal{C}$ is 
indeed an additive category. 

Now, do kernel exist? Suppose it exists and call it $\text{Ker}(f)$ for some
non-zero $f: m \to m$. Then there should be a morphism $\iota: \text{Ker}(f) \to
m$ such that $f \circ \iota = 0$ and $\iota$ being universal. However, all
morphisms $\text{Ker}(f) \to m$ for $\text{Ker}(f) \neq m$ are zero and $f \circ
\text{id}_m = 1 \neq 0$, so kernel does not exist. Similar logic can be applied
to show that cokernel does not exist.

\begin{problem}
    Let $X=\{(x,y,z) \in \mathbb{R}^3: x^2+y^2=z^3,z>0\}$, $Y=\mathbb{R}^2$ and
    let $f: X \to Y$ be the map $(x,y,z) \mapsto (x,y)$. Get $G :=
    Rf_*\mathbf{k}_X$. Calculate the stalk $G_0$ of $G$ at $0 \in \mathbb{R}^2$.
\end{problem}

\textbf{Solution.} Recall that:
\begin{equation*}
    (Rf_*\mathbf{k}_X)_0 = \varinjlim_{U \ni 0}H(f^{-1}(U), \mathbf{k}_X)
\end{equation*}

Now note that $f^{-1}(U) \cong \mathbb{S}^1$ for any open neighbourhood 
$U \ni 0$, thus:
\begin{equation*}
    (Rf_*\mathbf{k}_X)_0 = \varinjlim_{U \ni 0}H(\mathbb{S}^1, \mathbf{k}_X) = H(\mathbb{S}^1, \mathbf{k}_X) \cong H(\mathbb{S}^1, \mathbf{k}_{\mathbb{S}^1}) \cong \boxed{\mathbf{k} \oplus \mathbf{k}[-1]}.
\end{equation*}

\begin{problem}
    Let $X=\mathbb{R}^2$, $Y=\mathbb{R}$ and let $f: X \to Y$ be the projection 
    $(x,y) \mapsto x$.
    \begin{enumerate}
        \item[(a)] Let $Z = \{(x,y) \in X: y = x^2\}$. Compute $Rf_*\mathbf{k}_{XZ}$;
        \item[(b)] Let $S = \{(x,y) \in X: xy = 1\}$. Compute $Rf_*\mathbf{k}_{XS}$. 
    \end{enumerate}
\end{problem}

\textbf{Solution.} \textit{Part (a).} Similarly to the previous problem:
\begin{equation*}
    (Rf_*\mathbf{k}_{XZ})_x = \varinjlim_{U \ni x} H(f^{-1}(U), \mathbf{k}_{XZ}) \cong \varinjlim_{U \ni x} H(f^{-1}(U) \cap Z, \mathbf{k}_{X})
\end{equation*}

Now note that $f^{-1}(U)$ gives the long vertical stripe, thus $f^{-1}(U) \cap
Z$ gives the open connected segment on parabola $y=x^2$, which is homeomorphic
to $Y$. Therefore,
\begin{equation*}
    (Rf_*\mathbf{k}_{XZ})_x \cong H(\mathbb{R}, \mathbf{k}_{\mathbb{R}^2}) \cong \mathbf{k}
\end{equation*}

\textit{Part (b).} Previous formula applies:
\begin{equation*}
    (Rf_*\mathbf{k}_{XS})_x \cong \varinjlim_{U \ni x} H(f^{-1}(U) \cap S, \mathbf{k}_{X})
\end{equation*}

Now the expression depends on whether $x=0$ or not. If $x \neq 0$, we get 
the same result as before since we get an open connected segment on one of the 
branches of hyperbola:
\begin{equation*}
    (Rf_*\mathbf{k}_{XS})_x \cong \mathbf{k}, \; x \neq 0.
\end{equation*}

However, when $x=0$, we get that $f^{-1}(U) \cap S$ gives two branches of
hyperbola, both homeomorphic to $\mathbb{R}$. Therefore, $f^{-1}(U) \cap S \cong
\mathbb{R} \sqcup \mathbb{R}$ and so our cohomology:
\begin{equation*}
    (Rf_*\mathbf{k}_{XS})_0 \cong \mathbf{k}^2
\end{equation*}

Therefore, stalkwise, our resultant sheaf is given by:
\begin{equation*}
    (Rf_*\mathbf{k}_{XS})_x \cong \begin{cases}
        \mathbf{k}^2, & x = 0, \\
        \mathbf{k}, & x \neq 0.
    \end{cases}
\end{equation*}

\begin{problem}
    Let $X$ be a real manifold and let $S$ be a closed subset. Give a necessary 
    and sufficient condition in order that $\Gamma_S(C_X^{\infty}) \neq 0$.
\end{problem}

\textbf{Solution.} Recall that $\Gamma_S(C_X^{\infty}) = \{f \in C_X^{\infty}: S
\subseteq \text{supp}(f)\}$. We claim that necessary and sufficient condition is
$\text{int}(S) \neq \emptyset$. Indeed, if $\Gamma_S(C_X^{\infty}) \neq 0$, then
exists some point $s \in S$ such that $f(s) \neq 0$. Thus, there exists some
open neighbourhood $U$ of $s$ such that $f(U) \neq 0$, resulting in
$\text{int}(S) \neq \emptyset$. The reverse direction is obvious.

\begin{problem}
    Consider the map $f: X \to Y$ (where $X=\mathbb{R}^2$, $Y=\mathbb{R}^3$)
    such that $(x,y) \mapsto (x,y,x+y+1)$. Let $Z=\{(x,y,z) \in \mathbb{R}^3:
    x^2+y^2=z^2\}$. Compute $f^{-1}\mathbf{k}_{YZ}$.
\end{problem}

\textbf{Solution.} Recall that $f^{-1}\mathbf{k}_{Y,Z} \cong
\mathbf{k}_{Y,f^{-1}(Z)}$. Thus it suffices to compute $f^{-1}(Z)$.
Now suppose $(a,b) \in f^{-1}(Z)$, then $f(a,b) \in Z$ and thus:
\begin{equation*}
    a^2 + b^2 = (a+b+1)^2 \implies 2ab+2a+2b+1 = 0
\end{equation*}

Thus $f^{-1}(Z)$ is exactly the hyperbola $H := \{(x,y) \in \mathbb{R}^2:
2xy+2x+2y+1=0\}$ and so $f^{-1}\mathbf{k}_{Y,Z} \cong \mathbf{k}_{YH}$.

\begin{problem}
    Consider the map $f: X \to Y$ (where $X=\mathbb{R}^2$, $Y=\mathbb{R}$)
    given by $(x,y) \mapsto y$. Let $Z=\{(x,y) \in \mathbb{R}^2: y=x^2\}$.
    Describe the sheaf $Rf_*\mathbf{k}_{YZ}$.
\end{problem}

\textbf{Solution.} As we used previously,
\begin{equation*}
    (Rf_*\mathbf{k}_{XS})_x = \varinjlim_{U \ni x} H(f^{-1}(U) \cap S, \mathbf{k}_X).
\end{equation*}

Now, $f^{-1}(U)$ gives the horizontal stripe. Note that for $x=x^{-} < 0$, there
is no intersection with $S$, so $(Rf_*\mathbf{k}_{XS})_{x^{-}} \cong 0$. For
$x=0$, there is a single connected component, so $(Rf_*\mathbf{k}_{XS})_0 \cong
\mathbf{k}$. Finally, for $x=x^+>0$, we get disjoint union of two contractible 
sets, thus yielding $(Rf_*\mathbf{k}_{XS})_{x^+} \cong \mathbf{k}^2$. All in all:
\begin{equation*}
    (Rf_*\mathbf{k}_{XS})_{t} \cong \begin{cases}
        \mathbf{k}^2 & t>0, \\
        \mathbf{k} & t=0, \\
        0 & t<0.
    \end{cases} 
\end{equation*}

\begin{problem}
    Let $V$ be a real finite dimensional vector space and let $V^*$ be 
    its dual. One denotes by $q_1$ and $q_2$ the projections 
    $V \times V^* \to V$ and $V \times V^* \to V^*$, respectively.
    Consider the closed set
    \begin{equation*}
        P = \{(x,y) \in V \times V^*: \langle x, y \rangle \geq 0\}.
    \end{equation*}
    For $G \in \text{Mod}(\mathbf{k}_V)$, one defines the Fourier 
    transform as
    \begin{equation*}
        \mathcal{F}(G) = Rq_{2!}((q_1^{-1}G)_P).
    \end{equation*}
    Let $\gamma \subseteq V$ be a closed convex proper cone. Calculate
    $\mathcal{F}(\mathbf{k}_{V\gamma})$. 
\end{problem}

\textbf{Solution.} Gave up.

\begin{problem}
    Let $X$ be a topological space. One says that $F \in \text{Mod}(\mathbf{k}_X)$
    satisfies the unique continuation principle (UCP, for short) if for 
    any non-empty connected open set $U \subseteq X$ and any closed subset 
    $Z \subseteq U$ such that $Z \neq U$, one has $\Gamma_Z(U; F) = 0$.
    \begin{enumerate}
        \item[(i)] Prove that if $F$ satisfies UCP, then so does $F|_M$ and 
        $\Gamma_M(F)$ for any closed subset $M$ of $X$.
        \item[(ii)] Prove that if $X$ is open in $\mathbb{R}^n$ and $F$ is 
        soft, then $F$ does not satisfy the UCP.
        \item[(iii)] Give an example of a sheaf $F$ which satisfies the UCP principle
        but such that $R\Gamma_Z(U;F) \neq 0$ for some closed subset $Z$ with 
        $\text{Int}(Z) = \emptyset$.
        \item[(iv)] Prove that if $X$ is open in $\mathbb{R}^n$ and $F$ is soft,
        then $F$ does not satisfy the UCP.
    \end{enumerate}
\end{problem}

\begin{problem}
    Let $M \in \text{Mod}(\mathbf{k})$. Show that, in general, the functor
    $\text{Hom}_{\mathbf{k}}(M,\bullet)$ does not commute with directed
    inductive limits.
\end{problem}

\textbf{Solution.}

\begin{problem}
    Assume $\mathbf{k}$ is a field of characteristic $0$ and denote by $W_n$
    Weyl algebra in $n$ variables over $\mathbf{k}$. We set 
    $O_n = W_n/(W_n\partial_1 + \dots + W_n\partial_n)$ and recall that 
    $O_n \cong \mathbf{k}[x_1,\dots,x_n]$ as a left $W_n$-module. We also 
    set $B_n = W_n/(W_n\partial_1^2+\dots+W_n\partial_n^2)$.
    \begin{itemize}
        \item[(i)] First assume $n=1$ and compute $W_1/(W_1\partial_1^2)$ and 
        the groups $\text{Ext}_{W_1}^j(B_1,O_1)$.
        \item[(ii)] Prove that there exists an isomorphism of $W_n$-modules 
        of $B_n$ with a finite direct sums of copies of $O_1$.
        \item[(iii)] Calculate the groups $\text{Ext}^j_{W_n}(B_n, O_n)$.
    \end{itemize}
\end{problem}

\textbf{Solution.} \textit{Part (i).} So recall that
$W_1=\mathbf{k}[x,\partial]/(\partial x - x\partial - 1)$. Now,
\begin{equation*}
    W_1 / (W_1\partial^2) \cong \mathbf{k}[x,\partial]/(\partial x - x\partial - 1, W_1\partial^2),
\end{equation*}

which is a $\mathbf{k}$-vector space with basis $\{x^i, x^i\partial: i \geq
0\}$. We find the projective resolution of $B_1$:
\begin{equation*}
    0 \to W_1 \xrightarrow{\cdot \partial^2} W_1 \to B_1 \to 0
\end{equation*}

Now we apply $\text{Hom}_{W_1}(\bullet, O_1)$ functor to the projective resolution:
\begin{equation*}
    0 \to \text{Hom}_{W_1}(W_1, O_1) \xrightarrow{\cdot \partial^2} \text{Hom}_{W_1}(W_1, O_1) \to 0
\end{equation*}

But notice that $\text{Hom}_{W_1}(W_1, O_1) \cong O_1$, so
\begin{equation*}
    0 \to O_1 \xrightarrow{\cdot \partial^2} O_1 \to 0
\end{equation*}

Compute cohomologies (other are zeros):
\begin{equation*}
    \text{Ext}_{W_1}^0(B_1, O_1) \cong \text{Ker}(O_1 \xrightarrow{\cdot \partial^2} O_1), \quad \text{Ext}_{W_1}^1(B_1, O_1) \cong \text{Coker}(O_1 \xrightarrow{\cdot \partial^2} O_1)
\end{equation*}

Now, note that $\text{Ker}(\cdot)$ is exactly those polynomials $f \in O_1$ such
that $\partial^2f=0$, thus $\text{Ext}_{W_1}^0(B_1, O_1) \cong \mathbf{k} \oplus
\mathbf{k}x$. Trivially $\text{Ext}_{W_1}^1(B_1, O_1) \cong O_1$.

\textit{Part (ii).} Essentially we have to compute:
\begin{equation*}
    W_n / (W_n\partial_1^2 + \dots + W_n\partial_n^2) \cong \frac{\mathbf{k}[x_1,\dots,x_n,\partial_1,\dots,\partial_n]}{(x_i\partial_j - x_j\partial_i - \delta_{ij}, W_n\partial_1^2 + \dots + W_n\partial_n^2)}
\end{equation*}

Now, initially $W_n$ is a infinite-dimensional vector space with basis 
\begin{equation*}
    \{x_1^{\alpha_1}\dots x_n^{\alpha_n}\partial_1^{\beta_1}\dots \partial_n^{\beta_n}: \alpha_i, \beta_i > 0\}
\end{equation*}

And we ``kill'' all $\partial_i^{\alpha_i}$ for which $\alpha_i \geq 2$.
Therefore, $W_n / (W_n\partial_1^2 + \dots + W_n\partial_n^2)$ is a vector space
spanned by the basis $\{x_1^{\alpha_1}\dots x_n^{\alpha_n}\partial_1^{b_1}\dots
\partial_n^{b_n}: b_i \in \{0,1\}, \alpha_i \geq 0\}$. Now, define the map:
\begin{equation*}
    \sum_{b \in \{0,1\}^n} \partial_1^{b_1}\dots \partial_n^{b_n}f_{b_1,\dots,b_n}(x) \mapsto (f_{b_1,\dots,b_n})_{b \in \{0,1\}^n}
\end{equation*}

Thus, $B_n = W_n/(W_n\partial_1^2 + \dots + W_n\partial_n^2) \cong O_n^{\oplus 2^n}$.

\textit{Part (iii).} Note that we have:
\begin{equation*}
    \text{Ext}^j_{W_n}(B_n, O_n) \cong \text{Ext}^j_{W_n}(O_n^{\oplus 2^n}, O_n) \cong \prod_{2^n}\text{Ext}^j_{W_n}(O_n, O_n)
\end{equation*}

Apply Künneth formula:
\begin{equation*}
    \text{Ext}_{W_n}^*(O_n, O_n) \cong \bigotimes_{i=1}^n \text{Ext}_{W_1}^*(O_1, O_1).
\end{equation*}

Since $\text{Ext}_{W_1}^*(O_1, O_1) \cong \mathbf{k}[0]$, we obtain:
\begin{equation*}
    \text{Ext}^j_{W_n}(B_n, O_n) \cong \mathbf{k}^{2^n}
\end{equation*}

\textbf{Note.} General Künneth formula:
\begin{equation*}
    \text{Ext}_{R \otimes S}^n(M \otimes N, M' \otimes N') \cong \bigoplus_{p+q=n}(\text{Ext}_R^p(M,M') \otimes_{\mathbf{k}} \text{Ext}_S^q(N, N'))
\end{equation*}

\begin{problem}
    We assume that $\mathbf{k}$ is a field of characteristic $0$ and we denote 
    by $W_n$ the Weyl algebra in $n$ variables.
    \begin{enumerate}
        \item[(i)] Let $I=W_n\cdot x_1^2+\dots+W_n \cdot x_n^2$ be the left
        ideal generated by $(x_1^2,\dots,x_n^2)$. Show that a suitable 
        Koszul complex gives a finite free resolution of $W_n/I$.
        \item[(ii)] Calculate $\text{RHom}_{W_n}(W_n/I, W_n/I)$.
    \end{enumerate}
\end{problem}

\textbf{Solution.} \textit{Part (i).} We build Koszul complex
$K^{\bullet}(W_n/I, (x_1^2,\dots,x_n^2))$. We get:
\begin{equation*}
    0 \to M^{(0)} \xrightarrow{d^0} \dots \xrightarrow{d^{n-1}} M^{(n)} \to 0,
\end{equation*}

Let us see what each $M^{(j)}$ and $d^j$ stands for. Let $(e_1,\dots,e_n)$ be
the canonical basis of $\mathbf{k}^n$, denote $e_I = e_{i_1} \wedge \dots \wedge
e_{i_j}$ for $I = \{i_1,\dots,i_j\} \subseteq \{1,\dots,n\}$. Then, $M^{(j)} =
W_n \otimes \Lambda^j \mathbf{k}^n \cong W_n^{\oplus C_n^j}$. Now,
the differential map is defined as:
\begin{equation*}
    d^j(w \otimes e_I) = \sum_{i=1}^n wx_i^2 \otimes e_i \wedge e_I.
\end{equation*}

Note that elements $(x_1^2,\dots,x_n^2)$ form a regular sequence since each is 
a non-zero divisor of $W_n$ and $W_n$ is a domain. Now, apply the 
contravariant functor $\text{Hom}(\bullet, W_n/I)$:
\begin{equation*}
    0 \to \text{Hom}(M^{(n)}, W_n/I) \xrightarrow{d_0^*} \dots \xrightarrow{d_{n-1}^*} \text{Hom}(M^{(0)}, W_n/I) \to 0
\end{equation*}

Now note that $M^{(j)} \cong W_n^{\oplus \binom{n}{j}}$, so $\text{Hom}_{W_n}(M^{(n)},
W_n/I) \cong \text{Hom}_{W_n}(W_n^{\oplus \binom{n}{j}}, W_n/I) \cong (W_n/I)^{\oplus
\binom{n}{j}}$. Thus we get:
\begin{equation*}
    0 \to (W_n/I)^{\oplus\binom{n}{1}} \xrightarrow{d_0^*} \dots \xrightarrow{d_{n-1}^*} (W_n/I)^{\oplus\binom{n}{n}} \to 0,
\end{equation*}
where $d_j^*(f) = f \circ d^j$. One notices:
\begin{equation*}
    \text{RHom}_{W_n}(W_n/I, W_n/I) \cong \bigoplus_{j=0}^n (W_n/I)^{\binom{n}{j}}[-j]
\end{equation*}

\begin{problem}
    Let $P$ be a polynomial with coefficients in $\mathbb{C}$ of degree $m>1$,
    thus defining a continuous map from $X=\mathbb{C}$ to $Y=\mathbb{C}$. Set 
    $S = \{z \in X: P'(z) = 0\}$ and $Z=P(S) \subseteq Y$. Recall that 
    the restriction to $Y \setminus Z$ of the sheaf $P_*(\mathbf{k}_X)$ is 
    a locally constant sheaf of rank $m$.
    \begin{enumerate}
        \item[(i)] Let $P(z)=z^2$. Prove that the sheaf $P_*(\mathbf{k}_X)$
        is locally constant on $Y \setminus \{0\}$, is not constant on 
        $Y \setminus \{0\}$ and is not locally constant on $Y$.
        \item[(ii)] Let $P(z)=z^2(z-1)^2$. Find the points of $Y$ in a neighbourhood
        of which the sheaf $P_*(\mathbf{k}_X)$ is not locally constant.
    \end{enumerate}
\end{problem}

\textbf{Proof.} \textit{Part (i).} Recall $P_*\mathbf{k}_X(U) =
\mathbf{k}_X(P^t(U))$. So if $U$ does not contain critical points on $P$,
$P^t(U)$ splits into $m$ disjoint connected components and so
$P_*\mathbf{k}_X(U) \cong \mathbf{k}^m$. 

Apply this to $P(z)=z^2$. The critical point is $z=0$, so over $Y \setminus
\{0\}$, $P_*\mathbf{k}_X|_{Y \setminus \{0\}}(U) \cong \mathbf{k}^2$, thus
locally constant. Let us show that it is not constant. Consider
\begin{equation*}
    \Gamma(\mathbb{R}_{>0}, P_*\mathbf{k}_X) = \varinjlim_{U \subseteq \mathbb{R}_{>0}} \Gamma(P^{-1}(U), \mathbf{k}_X) \cong \mathbf{k},
\end{equation*}
which is not isomorphic to $\mathbf{k}^2$, as was expected. Now consider 
the whole $Y$. Let us compute the stalk at $0$:
\begin{equation*}
    (P_*\mathbf{k}_X)_0 = \varinjlim_{U \ni 0} \Gamma(P^{-1}(U), \mathbf{k}_X) \cong \mathbf{k}
\end{equation*}

Which differs from any stalk $w \neq 0$, where $(P_*\mathbf{k}_X)_w \cong
\mathbf{k}^2$. Thus, $P_*\mathbf{k}_X(U)$ is not locally constant over whole
$Y$.

\textit{Part (ii).} As was previously noted, $P_*(\mathbf{k}_X)|_{Y \setminus
Z}$ is locally constant of rank $4$ for $S=\{z \in X: P'(z) = 0\}$ and $Z=P(S)$. Compute the
derivative:
\begin{equation*}
    P'(z) = 2z(z-1)^2 + 2z^2(z-1) = 2z(z-1)(2z-1)
\end{equation*}

Thus $Z=P(\{0,1,\frac{1}{2}\})=\{0,1/16\}$.

\begin{problem}
    Let $V$ be an infinite dimensional $\mathbf{k}$-vector space. Denote by 
    $I$ the set of finite dimensional vector subspaces of $V$. For sake of 
    clarity, we denote by $V_i$ the vector space associated to $i \in I$.
    We order $I$ by the relation $i \leq j$ if $V_i \subseteq V_j$.
    \begin{enumerate}
        \item[(i)] Prove that $I$ is directed.
        \item[(ii)] Compute $\varinjlim_{i \in I} V/V_i$. 
    \end{enumerate}
\end{problem}

\textit{Part (i).} We need to show (i) $I$ is non-empty, (ii) for any $i$ and
$j$ in $I$, there exists $k \in I$ and morphisms $i \to k, j \to k$, and (iii)
for any parallel morphisms $f,g: i \rightrightarrows j$, there exists a morphism
$h: j \to k$ such that $h \circ f = h \circ g$.

(i) is obvious: any vector space has a finite dimensional subspace. For (ii),
assume $V_i$ and $V_j$ are two finite vector spaces. Then, there is always a
finite vector space $V_i + V_j$ with trivial maps $V_i \to V_i + V_j$ and $V_j
\to V_i + V_j$. Obviously, $V_i + V_j$ is finite dimensional, thus lies in $I$.
(iii) follows from $I$ having an inclusion order.

\textit{Part (ii).} By the definition of a colimit,
\begin{equation*}
    \varinjlim_{i \in I} V/V_i = \left(\bigsqcup_{i \in I}V/V_i\right) \Big/ \sim,
\end{equation*}

where $v+V_i \sim w+V_j$ iff exists $V_k \supseteq V_i,V_j$ such that
$v+V_k=w+V_k$ in $V/V_k$. But this obviously always holds, thus $\varinjlim_{i
\in I} V/V_i=0$.

\begin{problem}
    Let $C_c^{\infty}(\mathbb{R}^n)$ denotes the $\mathbb{C}$-vector 
    space of $\mathbb{C}$-valued $C^{\infty}$-functions on $\mathbb{R}^n$
    with compact support. 
    \begin{itemize}
        \item[(i)] Let $\theta \in C_c^{\infty}(\mathbb{R})$ with
        $\int_{\mathbb{R}}\theta(t)dt=1$. Prove that any $\varphi \in
        C_c^{\infty}(\mathbb{R})$ may be uniquely decomposed as $\varphi =
        \varphi_0 + \lambda\theta$ with $\varphi_0 \in
        C_c^{\infty}(\mathbb{R})$, $\lambda \in \mathbb{R}$, and
        $\int_{\mathbb{R}}\varphi_0(t)dt=0$. Also prove that there exists $\psi
        \in C_c^{\infty}(\mathbb{R})$ such that $\frac{\partial}{\partial
        x_1}\psi = \varphi_0$.
        \item[(ii)] Consider the operator $\frac{\partial}{\partial t}:
        C_c^{\infty}(\mathbb{R}) \to C_c^{\infty}(\mathbb{R})$. Prove that 
        its kernel is $0$ and its cokernel has a dimension $1$.
        \item[(iii)] We write $x=(x_1,x') \in \mathbb{R}^n$,
        $x'=(x_2,\dots,x_n)$. Prove that any $\varphi \in
        C_c^{\infty}(\mathbb{R}^n)$ may be uniquely decomposed as $\varphi(x) =
        \partial_{x_1}\psi + \lambda(x')\theta(x_1)$ with $\psi \in
        C_c^{\infty}(\mathbb{R}^n)$ and $\lambda(x') \in
        C_c^{\infty}(\mathbb{R}^{n-1})$.
        \item[(iv)] Consider the complex
        \begin{equation*}
            0 \to \Omega_c^0(\mathbb{R}^n) \xrightarrow{d} \dots \xrightarrow{d} \Omega_c^n(\mathbb{R}^n) \to 0,
        \end{equation*}
        where $\Omega_c^p(\mathbb{R}^n)$ denotes the space of $p$-differential
        forms on $\mathbb{R}^n$ with $C^{\infty}$-coefficients and compact
        support. Show that this complex is a Koszul complex and compute its
        cohomology.
    \end{itemize}
\end{problem}

\textit{Part (i).} By applying $\int_{\mathbb{R}}(\cdot)$ to both parts of the
equation we immediately get $\lambda=\int_{\mathbb{R}}\varphi(t)dt$. Define
$\varphi_0 := \varphi - \lambda\theta$. Its integral is obviously $0$ and it is
in $C_c^{\infty}(\mathbb{R})$. Uniqueness is obvious.

On the second part, we can pick $\psi(x_1) :=
\int_{-\infty}^{x_1}\varphi_0(t)dt$. It is clearly in $C^{\infty}(\mathbb{R})$,
so we shall only show that $\psi$ vanishes outside some $[-\rho,\rho]$. Since
$\varphi_0 \in C_c^{\infty}(\mathbb{R})$, it vanishes outside some
$[-\gamma,\gamma]$, thus $\psi$ vanishes over $(-\infty,-\gamma]$. Now, since
$\int_{\mathbb{R}}\varphi_0(t)dt=0$, thus for any $r>\gamma$, we have:
\begin{equation*}
    0 = \psi(\infty) = \psi(r) + \int_{r}^{\infty}\varphi_0(t)dt,
\end{equation*}
but the latter term is clearly $0$, thus $\psi(r)=0$. Thus, $\psi$ vanishes
outside $[-\gamma,\gamma]$.

\textit{Part (ii).} Kernel of this map is any $f \in C_c^{\infty}(\mathbb{R})$
such that $\partial_t f = 0$, meaning $f$ is a constant function. However,
constant functions are not in $C_c^{\infty}(\mathbb{R})$ (except for $0$), thus
the kernel is $0$.

Cokernel is, by definition, $C_c^{\infty}(\mathbb{R})/\text{im}(\partial_t)$. So
what exactly is the image of $\partial_t$? $g \in \text{im}(\partial_t)$ if
$g=\partial_t f$ for some $f \in C_c^{\infty}(\mathbb{R})$. Clearly if
$\int_{\mathbb{R}}g(t)dt=0$, then $g \in \text{im}(\partial_t)$ by \textit{part
(i)}. In general, it is easy to show that if $\int_{\mathbb{R}}g(t)dt \neq 0$,
then there are no functions $h \in C_c^{\infty}(\mathbb{R})$ such that
$\partial_t h = g$. Therefore, $\text{im}(\partial_t) = \{f \in
C_c^{\infty}(\mathbb{R}): \int_{\mathbb{R}}f(t)dt=0\}$ and thus
$C_c^{\infty}(\mathbb{R})/\text{im}(\partial_t)$ gives equivalence classes where
$f \sim g$ iff $\int_{\mathbb{R}}f(t)dt=\int_{\mathbb{R}}g(t)dt$. But notice
that the equivalence class is fully determined by the value of the integral
$\int_{\mathbb{R}}f(t)dt$. Thus, $C_c^{\infty}(\mathbb{R})/\text{im}(\partial_t)
\cong \mathbb{C}$ is of rank $1$.

\textit{Part (iii).} Apply $\int_{\mathbb{R}}(\bullet)(t,x')dt$ to both sides.
We get $\int_{\mathbb{R}}\varphi(t,x')dt = \lambda(x')$. Thus, define $\psi$
such that $\partial_{x_1}\psi = \varphi(x) - \lambda(x')\theta(x_1)$. Notice
that $\varphi(x) - \lambda(x')\theta(x_1) \in C_c^{\infty}(\mathbb{R}^n)$
automatically. Thus our proof boiled down to showing:

\textbf{Claim.} For any $f \in C_c^{\infty}(\mathbb{R}^n)$ such that
$\int_{\mathbb{R}}f(t,x')dt=0$, there exists a unique $g \in
C_c^{\infty}(\mathbb{R}^n)$ such that $f=\partial_{x_1}g$.

This is easily checked by induction for $g(x) = \int_{-\infty}^{x_1}f(t,x')dt$.

\textit{Part (iv).} Note that $\partial=(\partial_{x_1},\dots,\partial_{x_n})$
are trivially $\mathbb{C}$-linear endomorphisms of $C_{\infty}^c(\mathbb{R}^n)$
which commite with one another: $[\partial_{x_i},\partial_{x_j}] = 0$. Thus,
from $\Omega^p(\mathbb{R}^n)$ being a Koszul complex, we immediately get that
this is a Koszul complex $K^{\bullet}(C_c^{\infty}(X), \partial)$.

\begin{problem}
    Let $\mathcal{C}$ be an additive category.
    \begin{itemize}
        \item[(i)] Let $u: X_0 \to Y_0$ and $v: X_1 \to Y_1$ be two morphisms 
        in $\mathcal{C}$. Prove that if $u \oplus v: X_0 \oplus X_1 \to Y_0 \oplus Y_1$
        is an isomorphism, then both $u$ and $v$ are isomorphisms.
        \item[(ii)] Give an example of an isomorphism $w: X_0 \oplus X_1 \to Y_0 \oplus Y_1$
        for which $X_0$ and $Y_0$ are not isomorphic.
    \end{itemize}
\end{problem}

\textbf{Solution.} \textit{Part (i).} Since $u \oplus v = \begin{bmatrix}
    u & 0 \\ 0 & v
\end{bmatrix}$ is an isomorphism, these exists some matrix $\begin{bmatrix}
    a & b \\ c & d
\end{bmatrix}$ with $a,c: Y_0 \to X_0$ and $b,d: Y_1 \to X_1$ such that:
\begin{equation*}
    \begin{bmatrix}
    u & 0 \\ 0 & v
\end{bmatrix}\begin{bmatrix}
    a & b \\ c & d
\end{bmatrix} = \begin{bmatrix}
    \text{id}_{Y_0} & 0 \\
    0 & \text{id}_{Y_1}
\end{bmatrix}, \quad \begin{bmatrix}
    a & b \\ c & d
\end{bmatrix}\begin{bmatrix}
    u & 0 \\ 0 & v
\end{bmatrix} = \begin{bmatrix}
    \text{id}_{X_0} & 0 \\
    0 & \text{id}_{X_1}
\end{bmatrix}.
\end{equation*}

From here one deduces $b=c=0$ and:
\begin{equation*}
    u \circ a = \text{id}_{Y_0}, \quad a \circ u = \text{id}_{X_0}, \quad 
    v \circ d = \text{id}_{Y_1}, \quad d \circ v = \text{id}_{Y_0},
\end{equation*}
which naturally induce an isomorphism with $u^{-1}=a$, $v^{-1}=d$.

\textit{Part (ii).} Take $w: X \oplus Y \to Y \oplus X$ induced 
by matrix $\begin{bmatrix}
    0 & 1 \\
    1 & 0 \end{bmatrix}$.

\begin{problem}
    Let $\mathbf{k}$ be a field and denote by $\text{Mod}^f(\mathbf{k})$ the
    abelian category of finite dimensional $\mathbf{k}$-vector spaces. For
    $V^{\bullet} \in C^b(\text{Mod}^f(\mathbf{k}))$ one defines its 
    Euler-Poincaré index $\chi(V)$ by the formula:
    \begin{equation*}
        \chi(V) = \sum_i (-1)^i\dim H^i(V).
    \end{equation*}
    \begin{itemize}
        \item[(i)] Prove that if the complex $V$ is exact, then $\chi(V)=0$.
        \item[(ii)] Let $0 \to V' \to V \to V'' \to 0$ be an exact sequence in 
        $C^b(\text{Mod}^f(\mathbf{k}))$. Prove that $\chi(V)=\chi(V')+\chi(V'')$.
        \item[(iii)] Let $V,W \in C^b(\text{Mod}^f(\mathbf{k}))$. Prove that
        $\chi(V \otimes W) = \chi(V)\chi(W)$.
    \end{itemize}
\end{problem}

\textbf{Solution.} \textit{Part (i).} By the definition of exactness,
$\text{im}(d_V^{i-1}) \cong \text{ker}(d_V^{i})$, thus cohomology groups are
$H^i(V) = \text{ker}(d_V^i)/\text{im}(d_V^{i-1}) \cong 0$. Trivially, this
implies $\chi(V)=0$.

\textit{Part (ii).} A short exact sequence of complexes induces a long 
exact sequence in cohomology:
\begin{equation*}
    \dots \to H^i(V') \to H^i(V) \to H^i(V'') \xrightarrow{\delta^i} H^{i+1}(V') \to \dots
\end{equation*}
Set $a_i := \dim\text{Im}(H^i(V') \to H^i(V))$, $b_i :=
\dim\text{Im}(H^i(V) \to H^i(V''))$ and $c_i := \dim \text{Im}(\delta^i)$. Then:
\begin{gather*}
    \dim H^i(V') = a_{i-1} + c_{i-1},\\
    \dim H^i(V) = a_i + b_i, \\
    \dim H^i(V'') = b_i + c_i.
\end{gather*}

One then easily checks that $\chi(V) - \chi(V') - \chi(V'') = \sum_i
(-1)^i(a_{i-1}-a_i+c_{i-1}-c_i)=0$.

\textit{Part (iii).} We apply the Künneth formula:
\begin{equation*}
    H^n(V \otimes W) \cong \bigoplus_{p+q=n} H^p(V) \otimes H^q(W).
\end{equation*}

Note that $\dim(H^p(V) \otimes H^q(W)) = \dim H^p(V) \dim H^q(W)$, thus:
\begin{equation*}
    \dim H^n(V \otimes W) = \sum_{p+q=n}\dim H^p(V) \dim H^q(W)
\end{equation*}

Now compute $\chi(V \otimes W)$:
\begin{align*}
    \chi(V \otimes W) = \sum_i (-1)^i\dim H^i(V \otimes W) &= \sum_i (-1)^i\sum_{p+q=i}\dim H^p(V) \dim H^q(W)\\
    &= \sum_i (-1)^i\dim H^i(V) \cdot \sum_j (-1)^j\dim H^j(W) \\
    &= \chi(V)\chi(W).
\end{align*}

\begin{problem}
    Let $X=\mathbb{R}$ and let $F = \bigoplus_{n>0}\mathbf{k}_{1/n}$, where
    $\mathbf{k}_{1/n}$ is the skyscraper sheaf supported by the point $1/n$.
    Calculate the stalk $F_0$.
\end{problem}

\textbf{Solution.} Note:
\begin{equation*}
    \Gamma(U, F) = \Gamma\left(U,\bigoplus_{n>0}\mathbf{k}_{1/n}\right) = \bigoplus_{1/n \in U}\mathbf{k}
\end{equation*}

Thus $F_0 = \varinjlim_{0 \ni U}\Gamma(U, F) = \varinjlim_{0 \ni
U}\bigoplus_{1/n \in U}\mathbf{k}$. However, we can take $U$ arbitrarily small
to make corresponding sections $0$, thus $F_0=0$.

\begin{problem}
    Let $\mathcal{C}$ be a category with the property that $X,Y \in \mathcal{C}$
    are isomorphic iff $X=Y$. Let $\mathcal{C} \cong \mathcal{C}'$ be an equivalence 
    of categories. Does $\mathcal{C}'$ have the same property as $\mathcal{C}$?
\end{problem}

\textbf{Solution.} Let us gather all the facts:
\begin{enumerate}
    \item $f: X \to X$ is an isomorphism in $\mathcal{C}$ if
exists $g: X \to X$ such that $f\circ g = g\circ f = \text{id}_X$.
    \item $\text{Hom}_{\mathcal{C}}(X,Y) \cong \text{Hom}_{\mathcal{C}'}(F(X), F(Y))$.
    \item For each $Y \in \mathcal{C}'$ exists $X \in \mathcal{C}$ and $F(X)
    \cong Y$.
\end{enumerate}

Notice condition $3$: the inverse of the given $Y \in \mathcal{C}'$ might not be
unique. This way, while we have an isomorphism $X \to X$, we might still have
isomorphisms of form $X' \to Y'$ in $\mathcal{C}'$ for $X' \neq Y'$.

Suppose $\text{Ob}(\mathcal{C}) = \{\mathbf{k}^1, \mathbf{k}^2, \dots\}$ and
maps are linear maps. Take $\text{Ob}(\mathcal{C}')=\text{Vect}(\mathbf{k})$.
Take functor $F: \mathbf{k}^n \mapsto \mathbf{k}^n$ that assings to
$\mathbf{k}^n$ itself. Note that this is an equivalence: for any $V \in
\text{Vect}(\mathbf{k})$ there exists $W \in \mathcal{C}$ (namely,
$W=\mathbf{k}^{\dim V}$) such that $F(W) \cong V$. Isomomorphism of Hom's is
obvious.

Note that for example $\mathbf{k}[x]_{\leq 1} \cong \mathbf{k}^2$, but they are
not equal. 

\begin{problem}
    Let $F: \mathcal{C} \to \mathcal{C}'$ be an additive functor of abelian 
    categories.
    \begin{enumerate}
        \item[(i)] Prove that if $F$ is fully faithful, then it is conservative.
        \item[(ii)] Give an example of a functor which is conservative, but not 
        fully faithful.
        \item[(iii)] Give an example of a functor which is full but not faithful.
        \item[(iv)] Give an example of a functor which is faithful but not full.  
    \end{enumerate}
\end{problem}

\textbf{Solution.} \textit{Part (i).} Assume $F(f) \in \text{Hom}(X',Y')$ is an
isomorphism between $X',Y' \in \mathcal{C}'$. By essential surjectivity, there
exist $X,Y$ such that $F(X) \cong X'$, $F(Y) \cong Y'$. Moreover, since
$\text{Hom}_{\mathcal{C}}(X,Y) \cong \text{Hom}_{\mathcal{C}'}(X',Y')$ due to
fully faithfulness, there exists $g \in \text{Hom}_{\mathcal{C}'}(X',Y')$ such
that $F(f) \cong g$. We claim $g=f$.

\textit{Part (ii).} Forgetful functor $\mathbf{Ab} \to \mathbf{Set}$.

\textit{Part (iii).} Any map from non-trivial category to the category with 1 object.

\textit{Part (iv).} Forgetful functor $\textbf{Ring} \to \textbf{Ab}$.

\end{document}